% s11-abelian.tex —— §9 交换的情形（原文 p.42–p.46 上半）
% 底本：Laurent Fargues, Geometrization of the local Langlands correspondence: an overview, §9
% 译手：abelian
%
% 说明：原文以下划线标记常值群（\ul{E}^\times、\ul{T}(E)、\ul{\QQ}_p^\times）与 Hecke 叠
% （\ul{Hecke}），译文照样保留；\ol{\QQ}_\ell 在原文中反复出现，这里用一个局部宏简化。

\section{交换的情形}% 原文 9. The abelian case

\subsection{\texorpdfstring{$\mathrm{GL}_1$}{GL1} 情形}% 原文 9.1. The GL1-case

\leavevmode\srcpage{42}本节对 $p$-进域 $E$ 上的 $\mathrm{GL}_1$ 验证该猜想。此时按线丛的次数 $d$ 有分解
\[
\Bun = \coprod_{d \in \ZZ} \Bun^d .
\]
固定 $E$ 的一个一致化元 $\pi$。由此得到曲线上一个典范线丛 $\cO(1)$，它对应 $\pi^{-1} \in \mathrm{GL}_1(L)$。这一选择固定后，有同构
\[
[\Spa(\FF_p)/\ul{E}^\times] \xrightarrow{\sim} \Bun^d .
\]
记 $\mathcal{T}_d$ 为 $\Bun^d$ 上相应的万有 $\ul{\QQ}_p^\times$-torsor\yizhu{原文用下划线记号 $\ul{E}^\times$、$\ul{T}(E)$、$\ul{\QQ}_p^\times$ 表示相应的「常值群」，即把抽象群看作 pro-étale 层；Hecke 叠也写作 $\ul{Hecke}$。译文照原文保留这些下划线，不另立译名。}{A：原文 §9.1 p.42 与 §9.2 p.44 的页面图像；I：本工程 glossary 未收「常值群」条}。它是 $\cO(d)$ 与 $\Bun^d$ 上万有线丛之间同构构成的 torsor\yizhu{torsor 指群的主齐性空间（principal homogeneous space）。本译本体例：一律保留英文 \texttt{torsor}，首次出现处括注中文；与全库术语表 T1190（候选译名「主齐性空间」）的差异见 \texttt{\_meta/冲突台账.tsv}。}{I：本库 T1190「torsor=主齐性空间」（D 级候选）；本工程 glossary 已收「torsor」条}。令 $\mu(z) = z^k$，并记
\[
\ul{Hecke}^d_\mu = \vec{h}^{-1}\bigl(\Bun^d \times \Spa(E)^\diamond\bigr).
\]
Hecke 对应限制为
\[
\begin{tikzcd}[column sep=large]
& \ul{Hecke}^d_\mu \arrow[dl, "\overleftarrow{h}"'] \arrow[dr, "\vec{h}"] & \\
\Bun^{d+k} & & \Bun^d \times \Spa(E)^\diamond .
\end{tikzcd}
\]

固定 $\cO_E$ 上一个与 $\pi$ 相伴的 Lubin–Tate 群，记 $E(1)$ 为其 Tate 模，$E(k) = E(1)^{\otimes k}$。这给出 $\Spa(E)^\diamond$ 上一个秩 $1$ pro-étale 局部系统，从而给出 $\Spa(E)^\diamond$ 上一个 $\ul{E}^\times$-torsor $\mathcal{LT}_k$。我们用记号 $\wedge$ 表示对两个 $\ul{E}^\times$-torsor 作 $\times_{\ul{E}^\times}$。例如有
\[
\mathcal{LT}_k \wedge \mathcal{LT}_{k'} = \mathcal{LT}_{k+k'} .
\]

\begin{proposition}\label{prop:9.1}
\begin{enumerate}[label=(\arabic*),leftmargin=*]
\item 态射 $\vec{h}$ 是同构。
\item 存在 $\ul{E}^\times$-torsor 之间的自然同构
\[
\overleftarrow{h}^* \mathcal{T}_{d+k} \simeq \vec{h}^* \mathcal{T}_d \wedge \mathcal{LT}_{-k},
\]
其中 $\mathcal{LT}_{-k}$ 仍指它从 $\Spa(E)^\diamond$ 到 $\Bun \times \Spa(E)^\diamond$ 的拉回。
\end{enumerate}
\end{proposition}

\begin{proof}
第 (1) 点显然。第 (2) 点是 $p$-进 Hodge 理论基本正合列的推论。以 $\cO(1)$ 的幂作扭转，可设 $d \ge 0$（以 $\cO(1)$ 扭转给出同构 $\Bun^d \xrightarrow{\sim} \Bun^{d+1}$）。再取对偶修正，还可设 $k \ge 0$。

$\Spa(E)^\diamond$ 上有层正合列
\[
0 \to \mathbb{B}^{\varphi=\pi^d} \otimes_E E(k)^\diamond \to \mathbb{B}^{\varphi=\pi^{k+d}} \to \mathbb{B}^+_{dR}/\mathrm{Fil}^k \mathbb{B}_{dR} \to 0,
\]
其中回忆 $\mathbb{B}(S) = \Gamma(S, \cO_{Y_S})$。今设 $S$ 是完满空间，$f : (X_S)_{pro-\text{ét}} \to S_{pro-\text{ét}}$。对 $X_S$ 上次数为 $i$ 的线丛 $\mathcal{E}$，若 $\mathcal{T}$ 是 $\cO(i)$ 与 $\mathcal{T}$ 之间同构对应的 $\ul{E}^\times$-torsor\yizhuw{原文此处把两个对象都记作 $\mathcal{T}$（「$\mathcal{T}$ 是 $\cO(i)$ 与 $\mathcal{T}$ 之间同构构成的 torsor」），自指不通；按上下文应为 $\mathcal{T}$ 是 $\cO(i)$ 与 $\mathcal{E}$ 之间同构构成的 torsor。译文照字面保留，不改原文。}，则有
\[
Rf_* \mathcal{E} = \mathcal{T} \times_{\ul{E}^\times} \mathbb{B}^{\varphi=\pi^i} .
\]
由于 $\mathbb{B}^{\varphi=\pi^d} \otimes_E E(k)^\diamond = \mathcal{LT}_k \times_{\ul{E}^\times} \mathbb{B}^{\varphi=\pi^k}$，容易推出
\[
\vec{h}^* \mathcal{T}_d = \overleftarrow{h}^* \mathcal{T}_{d+k} \wedge \mathcal{LT}_k ,
\]
于是结论成立。
\end{proof}

\leavevmode\srcpage{43}现在设 $\varphi : W_E \to \Ql^\times$ 是连续特征。定义
\[
\chi : E^\times \xrightarrow[\sim]{\Art} W_E^{ab} \xrightarrow{\varphi} \Ql^\times,
\]
其中“Art”指 Artin 互反映射\yizhu{这里取算术 Artin 互反映射的规范化：一致化元 $\pi$ 对应 Frobenius。于是 $\chi=\varphi\circ\Art$ 把 $\pi$ 送到 $\varphi(\mathrm{Frob})$；由于 $\varphi$ 只要求连续，$\chi(\pi)$ 未必是 $\ell$-进单位，这正是下文 Weil 层不能下降为通常局部系统的原因。}{A：原文 §9.1 p.43；B：Serre, \emph{Local Fields}, 第 XI 章}，其规范化使得 Frobenius 对应一致化元。再定义 $\mathcal{G}$ 为 $\Bun$ 上秩 $1$ 的 $\Ql$-局部系统，使得
\[
\mathcal{G}|_{\Bun^d} = \mathcal{T}_d \times_{\ul{E}^\times,\chi} \Ql .
\]

\emph{这仍不是我们要找的局部系统。}事实上，按特征层性质\yizhu{本工程 glossary 把 character sheaf 与 eigensheaf 都译作「特征层」。此处 character sheaf property 指几何 Langlands 猜想中「逐点由特征标决定」的那条性质，与 Hecke 特征层（eigensheaf）的性质不是同一件事，读者勿混。}{I：本工程 glossary「character sheaf=特征层、eigensheaf=特征层」；A：原文 §9.1 p.43}，我们要找的 $\Bun_{\ol{\FF}_q}$ 上的 Weil 层并不总能下降为 $\Bun$ 上的层（在本节的场合，因为 $\chi(\pi)$ 可能不是 $\ell$-进单位）。于是定义 $\mathcal{F}_\varphi$ 为 $\Bun_{\ol{\FF}_q}$ 上的 Weil 层：去掉 Weil 下降数据后
\[
\mathcal{F}_\varphi = \mathcal{G}|_{\Bun_{\ol{\FF}_q}},
\]
而其 Frobenius 下降数据由 $\mathcal{G}$ 的典范下降数据乘以 $\Bun^d_{\ol{\FF}_p}$ 上的 $\chi(\pi)^{-d}$ 给出。

\begin{proposition}\label{prop:9.2}
局部系统 $\mathcal{F}_\varphi$ 满足猜想的假设。
\end{proposition}

\begin{proof}
先看 Hecke 性质。由命题~\ref{prop:9.1} 推出
\begin{equation}\tag{10}\label{eq:9.10}
\begin{aligned}
\vec{h}_!(\overleftarrow{h}^* \mathcal{G})|_{\Bun^d \times \Spa(E)^\diamond}
&= \mathcal{T}_d \times_{\ul{E}^\times} \mathcal{LT}_{-k} \times_{\ul{E}^\times,\chi} \Ql \\
&= \mathcal{G}|_{\Bun^d} \boxtimes (\chi \circ \chi_{\mathcal{LT}})^{-k},
\end{aligned}
\end{equation}
其中 $\chi_{\mathcal{LT}} : W_E^{ab} \to \cO_E^\times \subset E^\times$ 是 Lubin–Tate 特征。局部类域论在 Lubin–Tate 群挠点上的实现告诉我们，$\Art \circ \chi_{\mathcal{LT}}$ 满足
\begin{enumerate}[label=(\arabic*),leftmargin=*]
\item $\Art \circ \chi_{\mathcal{LT}}(\sigma) = 1$，若 $\sigma = \Art(\pi)$；
\item $\Art \circ \chi_{\mathcal{LT}}(\sigma) = \sigma^{-1}$，若 $\sigma$ 落在 $W_E$ 的惯性子群在 $W_E^{ab}$ 中的像里。
\end{enumerate}
由第 (2) 点推出
\[
\vec{h}_!(\overleftarrow{h}^* \mathcal{F}_\varphi)|_{\Bun^d_{\ol{\FF}_p} \times \Spa(\breve{E})^\diamond}
= \mathcal{F}_\varphi|_{\Bun^d_{\ol{\FF}_p}} \boxtimes r_\mu \circ \varphi|_{I_E},
\]
这里两边都去掉 Frobenius 下降数据。由式 (10) 与第 (1) 点推出：若给 $r_\mu \circ \varphi|_{I_E}$ 配上由 $r_\mu \circ \varphi$ 给出的 Weil 下降数据，这就是同构\yizhuw{原文此处作「this is an isomorphic」，缺少名词，当为「an isomorphism」之误；译文按「这是一个同构」读出。}。

特征层性质这样验证。取 $\delta \in E^\times$，$v_\pi(\delta) = -d$。于是在 $\Bun^d$ 上有两个 $\ul{E}^\times$-torsor：$\mathcal{T}_\delta$ 与 $\mathcal{T}_d = \mathcal{T}_{\pi^{-d}}$。考虑 $\ul{E}^\times$-齐性空间
\[
\Isom\bigl((L, \pi^{-d}\sigma), (L, \delta\sigma)\bigr)
= \{\, g \in L^\times \mid g\delta\sigma = \pi^{-d}\sigma g \,\}.
\]
对这样的 $g$，有 $g^\sigma = \delta\pi^d g$。由此推出
\[
\mathcal{T}_\delta = \mathcal{T}_d \wedge T_{\delta\pi^d},
\]
其中 $T_{\delta\pi^d}$ 是 $\FF_q$ 上的 $\ul{E}^\times$-torsor，其单值性由 $\Frob_q \mapsto \delta\pi^d$ 给出。特征层性质随即推出，因为我们在 $\delta = \pi^{-d}$ 处“强行”要求了它。
\end{proof}

\subsection{环面情形}% 原文 9.2. The case of tori

设 $T$ 是 $E$ 上的环面，$\Bun$ 为对应的叠。按 Kottwitz\yizhu{$\kappa$ 是 Kottwitz 给出的 $\sigma$-共轭类集 $\mathrm{B}(T)$ 与 $X_*(T)_\Gamma$ 之间的典范同构；下文按 $\kappa$ 的像分解 $\Bun$，$\Bun^{[\mu]}$ 即 $\kappa=[-\mu]$ 的那一块。}{A：原文 §9.2 p.43；B：Kottwitz, \emph{Isocrystals with additional structure}（原文文献 [18]）}，
\[
\kappa : \mathrm{B}(T) \xrightarrow{\sim} X_*(T)_\Gamma .
\]
有分解
\[
\Bun = \coprod_{[\mu] \in X_*(T)_\Gamma} \Bun^{[\mu]},
\]
其中记 $\Bun^{[\mu]} = \{\kappa = [-\mu]\}$（我们采用与猜想中相反的符号约定，如同 $\mathrm{GL}_1$ 情形；向量丛的次数取为 $\kappa$ 的相反数，也许我们该在整个猜想中都这样做\srcpageinl{44}）。

对 $\mu' \in X_*(T)$，记 $\bar\mu' \in X_*(T)/\Gamma$ 与 $[\mu'] \in X_*(T)_\Gamma$ 为其类。给定 $\mu$ 与 $\mu'$，有 Hecke 对应
\[
\begin{tikzcd}[column sep=large]
& \ul{Hecke}^{\bar\mu'} \arrow[dl, "\overleftarrow{h}"'] \arrow[dr, "\vec{h}"] & \\
\Bun^{[\mu+\mu']} & & \Bun^{[\mu]} \times \Spa(E)^\diamond .
\end{tikzcd}
\]

容易验证下述命题。

\begin{proposition}\label{prop:9.3}
若 $E_{\mu'}|E$ 满足 $\mathrm{Stab}_\Gamma(\mu') = \Gamma_{E_{\mu'}}$，则有等同 $\ul{Hecke}^{\bar\mu'} = \Bun \times \Spa(E_{\mu'})^\diamond$，其中 $\vec{h}$ 由 $\mathrm{Id} \times \mathrm{can}_{E_{\mu'}|E}$ 给出。
\end{proposition}

设 $\varphi : W_E \to {}^L T$ 为 L-参数。按局部类域论，$H^1(W_E, \hat{T}) = \Hom(T(E), \Ql)$。于是有特征
\[
\chi : T(E) \to \Ql^\times .
\]
对 $b \in T(L)$ 有等同
\[
x_b : [\Spa(\ol{\FF}_q)/\ul{T}(E)] \xrightarrow{\sim} \Bun^{-\kappa(b)}_{\ol{\FF}_q} .
\]
这给出 $\Bun^{-\kappa(b)}_{\ol{\FF}_q}$ 上的 $T(E)$-torsor $\mathcal{T}_b$。有 $\Frob^* x_b = x_{b^\sigma}$ 与 $\Frob^* \mathcal{T}_b = \mathcal{T}_{b^\sigma}$。$\sigma$-共轭 $b^\sigma = b^{-1}.b.b^\sigma$ 给出 $x_b$ 与 $x_{b^\sigma}$ 之间的自然变换。于是得到 Weil 下降数据
\begin{equation}\tag{11}\label{eq:9.11}
\Weil_b : \Frob^* \mathcal{T}_b \xrightarrow{\sim} \mathcal{T}_b .
\end{equation}

\begin{example}\label{ex:9.4}
若 $b \in T(E_n)$，则 $\Frob^{n*} x_b = x_b$。这给出 $\Frob^{n*} \mathcal{T}_b$ 与 $\mathcal{T}_b$ 之间的一个等同，也就是说 $\mathcal{T}_b$ 定义在 $\Bun^{-\kappa([b])}_{\ol{\FF}_{q^n}}$ 上。下降数据 (11) 的第 $n$ 次迭代就是这个等同乘以 $N_{E_n/E}(b)$。例如，若 $b \in T(E)$，则 $\mathcal{T}_b$ 上的这个 Weil 下降数据除非 $b$ 是单位根，否则绝不有效。
\end{example}

若 $b'$ 与 $b$ 是 $\sigma$-共轭的，令
\[
T_{b,b'} = \{\, g \in T(L) \mid gbg^{-\sigma} = b' \,\}.
\]
这是 $T(E)$ 下的齐性空间，我们把它看作 $\Spa(\ol{\FF}_q)$ 上的 $T(E)$-torsor。有 $\Bun^{-\kappa([b])}_{\ol{\FF}_q}$ 上 torsor 之间的等同
\[
\mathcal{T}_b \wedge T_{b,b'} = \mathcal{T}_{b'} .
\]
固定 $T_{b,b'}$ 的一个平凡化，即取 $g \in T(L)$ 使 $gbg^{-\sigma} = b'$。这给出 Weil 层的同构
\[
g : \mathcal{T}_b \xrightarrow{\sim} \mathcal{T}_{b'},
\]
意思是 $g \circ \Weil_b = \Weil_{b'} \circ \Frob^* g$，其中 $\Frob^* g = g^\sigma$。

设 $\varphi : W_E \to {}^L T$ 为 Langlands 参数，$\chi : T(E) \to \Ql^\times$ 为类域论给出的相应特征，$H^1(W_E, \hat{T}) = \Hom(T(E), \Ql^\times)$。我们按公式
\[
\mathcal{F}|_{\Bun^{[\mu]}_{\ol{\FF}_q}} = \mathcal{T}_b \times_{\ul{T}(E),\chi} \Ql
\]
定义 $\Bun_{\ol{\FF}_q}$ 上的秩一 Weil 局部系统 $\mathcal{F}$，其中取某个 $b \in T(L)$ 使 $\kappa([b]) = -[\mu]$。$b$ 的另一种取法不改变这个 Weil 局部系统的同构类（而猜想中要求的所有性质只依赖于\srcpageinl{45}这个 Weil 反常层\yizhu{原文此处说「Weil 反常层」，但 $\mathcal{F}$ 按定义是秩一的 Weil 局部系统；作者在 §9 里把「局部系统」与「Weil 反常层」混用，译文照原文。}{A：原文 §9.2 p.44–p.45；I：本工程 glossary「perverse sheaf=反常层」}的同构类，不依赖于 Weil 反常层本身）。

我们现在要证明 $\mathcal{F}$ 满足猜想的假设。特征层性质是直接的。唯一要证的是 Hecke 性质。取 $\mu, \mu' \in X_*(T)$，$E_{\mu'}$ 如前为 $\mu'$ 的定义域。令
\[
\delta = N_{E_{\mu'}/E_n}(\mu'(\pi_{\mu'})) \in T(E_n),
\]
其中 $E_n|E$ 是 $E_{\mu'}$ 中的最大非分歧扩张。有 $\kappa(\delta) = [\mu']$，并且 $(\delta, \mu')$ 是可容许对。这定义晶体表示\yizhuw{原文作「cristalline representation」（法语拼法），英文标准拼法为 crystalline；译文统一作「晶体表示」。}
\[
\rho_{\delta,\mu'} : \Gamma^{ab}_{E_{\mu'}} \to T(E).
\]
记
\[
N_{\mu'} : E^\times_{\mu'} \xrightarrow{\mu'} T(E_{\mu'}) \xrightarrow{N_{E_{\mu'}/E}} T(E).
\]
与 Artin 互反映射 $\Art : E^\times_{\mu'} \to \Gamma^{ab}_{E_{\mu'}}$ 复合后，它满足
\begin{itemize}[leftmargin=*]
\item $\rho_{\delta,\mu'} \circ \Art(x) = N_{\mu'}(x)^{-1}$，$x \in \cO^\times_{E_{\mu'}}$；
\item $\rho_{\delta,\mu'} \circ \Art(\pi_{\mu'}) = 1$。
\end{itemize}
我们记 $T^{cris}_{\delta,\mu'}$ 为 $\Spa(E_{\mu'})^\diamond$ 上相应的 $\ul{T}(E)$-torsor。

\begin{proposition}\label{prop:9.5}
对 $b \in T(L)$，$\ul{Hecke}^{\bar\mu'}_{\ol{\FF}_q}$ 上有 Weil $\ul{T}(E)$-torsor 的同构
\[
\overleftarrow{h}^* \mathcal{T}_{b\delta} \simeq \vec{h}^* \mathcal{T}_b \wedge T^{cris}_{\delta,\mu'} .
\]
\end{proposition}

\begin{proof}
取 $(V,\rho) \in \Rep_E(T)$，即 $T$ 的一个线性表示。注意 $\mathbb{C}_p = \widehat{\ol{E}}$。晶体表示 $\rho \circ \rho_{\delta,\mu'}$ 对应滤过等晶体 $(V_{E_n}, \rho(\delta)\sigma, \rho \circ \mu')$。这意味着在 $\infty \in |X_{\mathbb{C}_p^\flat}|$ 处有向量丛的 $\Gamma_{E_{\mu'}}$-等变修正（见 [10] 第 10 章）
\[
(V \otimes_E \cO_{X_{\mathbb{C}_p^\flat}})|_{X_{\mathbb{C}_p^\flat} \setminus \infty}
\xrightarrow{\sim}
\mathcal{E}(V_{E_n}, \rho(\delta)\sigma)|_{X_{\mathbb{C}_p^\flat} \setminus \infty},
\]
其中左端 $\Gamma_{E_{\mu'}}$ 的作用由 $\rho \circ \rho_{\delta,\mu'}$ 给出。这个修正的类型由 $\rho \circ \mu'$ 给出。以 $\rho_* \mathcal{E}_b = \mathcal{E}(V_L, \rho(b)\sigma)$ 扭转，得到 $\Gamma$-等变修正
\[
\rho_*(\mathcal{E}_b \wedge T^{cris}_{\delta,\mu'})|_{X_{\mathbb{C}_p^\flat} \setminus \infty}
\xrightarrow{\sim}
\rho_*(\mathcal{E}_{b\delta})|_{X_{\mathbb{C}_p^\flat} \setminus \infty}.
\]
经 $(X_{\mathbb{C}_p^\flat})_{pro-\text{ét}} \to \Spa(\mathbb{C}_p^\flat)_{pro-\text{ét}}$ 投影，计入 $\Gamma$-作用并让 $\rho$ 变动，即得 9.1 中的结果。
\end{proof}

Hecke 性质随后这样证明。余特征 $\mu'$ 定义\yizhuw{原文作「The chocaracter」，当为 cocharacter 之误；glossary 定 cocharacter 为「余特征」，译文据此。}
\[
\hat\mu' \in X^*(\hat{T})^{\Gamma_{E_{\mu'}}},
\]
它定义特征
\[
{}^L\mu' : {}^L T_{E_{\mu'}} \to \Ql^\times,
\]
其中 ${}^L T_{E_{\mu'}} = \hat{T} \rtimes \Gamma_{E_{\mu'}}$，且 ${}^L\mu'$ 在 $\Gamma_{E_{\mu'}}$-因子上平凡。于是有
\[
r_{\mu'} = \Ind_{E_{\mu'}|E} {}^L\mu' .
\]
局部类域论的同构
\[
H^1(W_E, \hat{T}) = \Hom(T(E), \Ql^\times)
\]
与基变换相容。这意味着 $\varphi|_{W_{E_{\mu'}}}$ 对应
\[
\chi \circ N_{E_{\mu'}/E} : T(E_{\mu'}) \to \Ql^\times .
\]

\leavevmode\srcpage{46}由此推出
\[
\chi \circ N_{E_{\mu'}/E} : E^\times_{\mu'} \xrightarrow[\sim]{\Art} W^{ab}_{E_{\mu'}} \xrightarrow{{}^L\hat\mu' \circ \varphi_{W_{E_{\mu'}}}} \Ql^\times .
\]
Hecke 性质由这一点以及前面关于 $\rho_{\delta,\mu'} \circ \Art$ 的公式推出。
